% \documentclass{article}
% \usepackage[utf8]{inputenc}

% \begin{document}

% \centerline{\bf DUNE Acknowledgement}
% \centerline{\bf January 2022}
% \bigskip

% Below is a standard paragraph to be included in DUNE publications. 
% The sentences acknowledging Fermilab and funding agencies should always appear. 
% The sentences relevant to ProtoDUNE would appear only in ProtoDUNE papers. 
% The sentence acknowledging NERSC should appear only in papers that used NERSC resources.
% \bigskip \hrule \bigskip

%
% Acknowledgement of CERN for ProtoDUNE construction and operations
%
The ProtoDUNE-SP detector was constructed and operated on the CERN Neutrino Platform.
We gratefully acknowledge the support of the CERN management, and the
CERN EP, BE, TE, EN and IT Departments for NP04/Proto\-DUNE-SP.
%
% Standard piece for all FNAL-based experiments
%
This document was prepared by the DUNE collaboration using the
resources of the Fermi National Accelerator Laboratory 
(Fermilab), a U.S. Department of Energy, Office of Science, 
HEP User Facility. Fermilab is managed by Fermi Research Alliance, 
LLC (FRA), acting under Contract No. DE-AC02-07CH11359.
%
% Funding agencies, alphabetical by country, then alphabetical by agency name
%
This work was supported by
CNPq,
FAPERJ,
FAPEG and 
FAPESP,                         Brazil;
CFI, 
IPP and 
NSERC,                          Canada;
CERN;
M\v{S}MT,                       Czech Republic;
ERDF, 
H2020-EU and 
MSCA,                           European Union;
CNRS/IN2P3 and
CEA,                            France;
INFN,                           Italy;
FCT,                            Portugal;
NRF,                            South Korea;
CAM, 
Fundaci\'{o}n ``La Caixa'',
Junta de Andaluc\'ia-FEDER,
MICINN, and
Xunta de Galicia,               Spain;
SERI and 
SNSF,                           Switzerland;
T\"UB\.ITAK,                    Turkey;
The Royal Society and 
UKRI/STFC,                      United Kingdom;
DOE and 
NSF,                            United States of America.
%
% Acknowledgement of NERSC if those resources were used
%
This research used resources of the 
National Energy Research Scientific Computing Center (NERSC), 
a U.S. Department of Energy Office of Science User Facility 
operated under Contract No. DE-AC02-05CH11231.
%
%\end{document}
